\subsection{紧线性映射}

定义 1. --- 设 E 是拓扑向量空间，F 是分离的拓扑向量空间。从 E 到 F 的线性映射 u 称为紧的，如果存在 E 中 0 的一个邻域 V，使得 \(u(V)\) 是 F 的一个相对紧子集。

我们用 \(\mathcal{L}^{\mathrm{c}}(\mathrm{E};\mathrm{F})\) 表示从 E 到 F 的紧线性映射所成的集；我们也用 \(\mathcal{L}^{\mathrm{c}}(\mathrm{F})\) 表示集 \(\mathcal{L}^{\mathrm{c}}(\mathrm{F};\mathrm{F})\)。

注记. --- 1) E 中 0 的一个邻域吸收 E 的每个有界子集；因此，E 的有界子集在从 E 到 F 的紧线性映射下的像是 F 的一个相对紧子集。

\begin{enumerate}
\def\labelenumi{\arabic{enumi})}
\setcounter{enumi}{1}
\item
  设 E 是半赋范空间，B 是 E 的单位球。为使从 E 到一个分离拓扑向量空间 F 的线性映射 u 是紧的，必须且只需 \(u(\mathrm{B})\) 是 F 的一个相对紧子集。
\item
  设 F 是分离的拓扑向量空间。为使 F 的向量子空间 E 是有限维的，必须且只需 E 到 F 中的典范单射是紧的（EVT, I, p. 15, th. 3）。
\item
  对域 R 上的每个拓扑向量空间 E，我们用 \(\mathrm{E}_{(\mathbf{C})}\) 表示 E 的复化拓扑向量空间（EVT, II, p. 65）。我们把 E 等同于 \(\mathrm{E}_{(\mathbf{C})}\) 的一个实拓扑向量子空间（经由映射 \(x \mapsto 1 \otimes x\)）；于是 \(\mathrm{E}_{(\mathbf{C})}\) 的底层实拓扑向量空间是 E 与 iE 的拓扑直和。
\end{enumerate}

设 \(u\) 是从 E 到 \(\mathbf{R}\) 上的一个分离拓扑向量空间 F 的线性映射。用 \(u_{(\mathbf{C})}\) 表示 \(\mathbf{C}\)-线性映射（它从 \(\mathrm{E}_{(\mathbf{C})}\) 到 \(\mathrm{F}_{(\mathbf{C})}\) 中，并且拓广 \(u\)）。形如 \(V + iV\) 的集合（其中 \(V\) 是 E 中 0 的一个邻域）构成 \(\mathrm{E}_{(\mathbf{C})}\) 中 0 的一个基本邻域系。为使集合 \(u_{(\mathbf{C})}(V + iV) = u(V) + iu(V)\) 在 \(\mathrm{F}_{(\mathbf{C})}\) 中相对紧，必须且只需 \(u(V)\) 在 F 中相对紧。因此，为使 \(\mathbf{C}\)-线性映射 \(u_{(\mathbf{C})}\) 是紧的，必须且只需 \(\mathbf{R}\)-线性映射 \(u\) 是紧的。

出于同样的思路，当 \(v\) 是从 \(\mathrm{E}_{(\mathbf{C})}\) 到一个 C-拓扑向量空间 G 的连续 C-线性映射时，为使 \(v\) 是紧的，必须且只需它在 E 上的限制是紧的。

\begin{enumerate}
\def\labelenumi{\arabic{enumi})}
\setcounter{enumi}{4}
\item
  设 E 是拓扑向量空间，F 是分离的拓扑向量空间。设 \(F_{1}\) 是 F 的一个闭向量子空间，j 是 \(F_{1}\) 到 F 中的典范单射。为使从 E 到 \(F_{1}\) 的线性映射 u 是紧的，必须且只需从 E 到 F 的线性映射 \(j \circ u\) 是紧的。
\item
  设 \((\mathrm{G}_{i})_{i\in\mathrm{I}}\) 是一族分离的拓扑向量空间，且对每个 \(i\in I\)，设 \(u_{i}\) 是从拓扑向量空间 E 到 \(G_{i}\) 的一个紧线性映射。若集合 I 有限，或者 E 是半赋范空间，则线性映射 \(x\mapsto(u_{i}(x))_{i\in\mathrm{I}}\)（从 E 到 \(\prod_{i\in I}G_{i}\) 中）是紧的（注记 2 及 TG, I, p. 63, Th. 3）。
\end{enumerate}

命题 1. --- 设 E 是拓扑向量空间，F 是分离的拓扑向量空间。

\begin{enumerate}
\def\labelenumi{\alph{enumi})}
\item
  从 E 到 F 的每个紧线性映射都是连续的。
\item
  从 E 到 F 的每个有限秩连续线性映射都是紧的；
\item
  从 E 到 F 的紧线性映射所成的集 \(\mathcal{L}^{\mathrm{c}}(\mathrm{E};\mathrm{F})\) 是 \(\mathcal{L}(\mathrm{E};\mathrm{F})\) 的一个向量子空间。
\end{enumerate}

我们来证明 a)。设 \(v: E \rightarrow F\) 是紧线性映射，U 是 F 中 0 的一个邻域。取 E 中 0 的一个邻域 V，使得 \(v(V)\) 在 F 中相对紧。于是 \(v(V)\) 有界，因此存在实数 \(\lambda > 0\)，使得集合 \(v(\lambda V)\) 包含于 U；从而 v 连续。

设 \(u \colon \mathrm{E} \to \mathrm{F}\) 是有限秩连续线性映射。由于 \(u(\mathrm{E})\)（即 \(u\) 的像）是分离的且是有限维的，故存在 \(u(\mathrm{E})\) 中 0 的一个紧邻域 A。集合 \(\mathrm{V} = u^{-1}(\mathrm{A})\) 是 E 中 0 的一个邻域，并且 \(u(\mathrm{V}) \subset \mathrm{A}\)，所以线性映射 \(u\) 是紧的。这就证明了断言 \(b\)）。

设 \(u_{1}\) 和 \(u_{2}\) 是从 E 到 F 的紧线性映射。取 E 中 0 的邻域 \(V_{1}\) 和 \(V_{2}\)，使得集合 \(u_{1}(V_{1})\) 与 \(u_{2}(V_{2})\) 都在 F 中相对紧。令 \(V = V_{1} \cap V_{2}\)，\(u = u_{1} + u_{2}\)。于是 \(u(V)\) 包含于 \(u_{1}(V_{1}) + u_{2}(V_{2})\)，从而在 F 中相对紧，故线性映射 u 是紧的。\(\mathcal{L}^{c}(E;F)\) 在位似变换下稳定是显然的。这就证明了 \(\mathcal{L}^{c}(E;F)\) 是 \(\mathcal{L}(E;F)\) 的一个向量子空间，从而得到断言 c)。

命题 2. --- 设 E 是半赋范空间，F 是分离的拓扑向量空间。假设 F 的每个有界闭子集都是完备的。那么向量子空间 \(\mathcal{L}^{\mathrm{c}}(\mathrm{E};\mathrm{F})\)（\(\mathcal{L}(\mathrm{E};\mathrm{F})\) 的）对有界收敛拓扑是闭的。

设 \(u\) 是 \(\mathcal{L}(\mathrm{E};\mathrm{F})\) 中对有界收敛拓扑附着于 \(\mathcal{L}^{\mathrm{c}}(\mathrm{E};\mathrm{F})\) 的一个元素。用 B 表示 E 的单位球，V 表示 F 中 0 的一个邻域。取 F 中 0 的一个邻域 W，使得 \(\mathrm{W} + \mathrm{W} \subset \mathrm{V}\)。按假设，存在元素 \(v\)（它属于 \(\mathcal{L}^{\mathrm{c}}(\mathrm{E};\mathrm{F})\)），使得 \((u - v)(\mathrm{B}) \subset \mathrm{W}\)。集合 \(v(\mathrm{B})\) 在 F 中相对紧，故存在 F 的一个有限子集 M，使得 \(v(\mathrm{B}) \subset \mathrm{M} + \mathrm{W}\)，从而 \(u(\mathrm{B}) \subset \mathrm{M} + \mathrm{V}\)。由于这对每个 V 都成立，集合 \(u(\mathrm{B})\) 在 F 中预紧，因而其闭包也预紧（TG, II, prop. 1, p. 30）。后者是闭的且有界，故由对 F 所作的假设，它是完备的，从而是紧的（TG, II, cor., p. 30）。由此推出 \(u(\mathrm{B})\) 相对紧。这就证明了 \(u\) 属于 \(\mathcal{L}^{\mathrm{c}}(\mathrm{E};\mathrm{F})\)。

若 F 是局部凸的，则命题的假设意味着 F 是拟完备的（EVT, III, p. 8, def. 6）。

推论. --- 设 E 和 F 是巴拿赫空间。集合 \(\mathcal{L}^{\mathrm{c}}(\mathrm{E};\mathrm{F})\) 是巴拿赫空间 \(\mathcal{L}(\mathrm{E};\mathrm{F})\) 的一个闭向量子空间。它包含从 E 到 F 的有限秩连续线性映射所成的空间 \(\mathcal{L}^{\mathrm{f}}(\mathrm{E};\mathrm{F})\) 的闭包。

这由命题 1 和命题 2 得出。

\pandocbounded{\includegraphics[keepaspectratio]{images/part001_746f5095.jpg}}

存在巴拿赫空间 E 和 F，使得 \(\mathcal{L}^{\mathrm{f}}(\mathrm{E};\mathrm{F})\) 在 \(\mathcal{L}^{\mathrm{c}}(\mathrm{E};\mathrm{F})\) 中不稠密（参见 III, p. 16, 注记 6 及 III, p. 19, 定理 4, b)）。

命题 3. --- 设 \(\mathrm{E}_{1}\)、E、F、\(\mathrm{F}_{1}\) 是拓扑向量空间，其中 F 和 \(\mathrm{F}_{1}\) 假设为分离的。设 v: \(\mathrm{E}_{1}\) → E、u: E → F、w: F → \(\mathrm{F}_{1}\) 是线性映射。若 v 和 w 连续而 u 紧，则 w ∘ u ∘ v 紧。

按假设，存在 E 中 0 的一个邻域 V，使得 \(u(\mathrm{V})\) 在 F 中相对紧。令 \(\mathrm{U} = v^{-1}(\mathrm{V})\)。于是 U 是 \(\mathrm{E}_{1}\) 中 0 的一个邻域，且它在 \(w \circ u \circ v\) 下的像包含于 \(w(u(\mathrm{V}))\)，从而在 \(\mathrm{F}_{1}\) 中相对紧。因此，线性映射 \(w \circ u \circ v\) 是紧的。

设 E 是分离的拓扑向量空间；由命题 1 和命题 3，\(\mathcal{L}^{\mathrm{c}}(\mathrm{E})\) 是代数 \(\mathcal{L}(\mathrm{E})\) 的一个双边理想。当 E 是巴拿赫空间时，\(\mathcal{L}(\mathrm{E})\) 是巴拿赫代数，而 \(\mathcal{L}^{\mathrm{c}}(\mathrm{E})\) 是 \(\mathcal{L}(\mathrm{E})\) 的一个闭双边理想（命题 2 的推论）。若 E 是无穷维的，则它是一个真理想（参见 p. 2, 注记 3）。

推论. — 设 E 是希尔伯特空间。空间 \(\mathcal{L}^{\mathrm{c}}(\mathrm{E})\) 是 \(\mathcal{L}(\mathrm{E})\) 的一个闭的自伴双边理想。特别地，它是一个星代数。

事实上，空间 \(\mathcal{L}^{c}(E)\) 是 \(\mathcal{L}(E)\) 的一个闭双边理想，从而是自伴的（I, p. 122, 引理 15）。

命题 4. — 设 E 是拓扑向量空间，\(\widehat{\mathrm{E}}\) 是它的分离完备化，\(j\) 是 E 到 \(\widehat{\mathrm{E}}\) 中的典范映射。设 \(u\) 是从 E 到一个分离拓扑向量空间 F 的紧线性映射。则存在唯一的紧线性映射 \(v\)（从 \(\widehat{\mathrm{E}}\) 到 F 中），使得 \(u = v \circ j\)。

把 F 等同于 \(\widehat{\mathrm{F}}\) 的一个拓扑向量子空间，并用 \(\widehat{u}\colon \widehat{\mathrm{E}}\to \widehat{\mathrm{F}}\) 表示使 \(\widehat{u}\circ j\) 在 E 上与 \(u\) 重合的唯一连续线性映射。由于 \(u\) 紧，存在 E 中 0 的一个邻域 V 和 F 的一个紧子集 A，使得 \(u(\mathrm{V})\subset \mathrm{A}\)。我们有 \(\widehat{u}(j(\mathrm{V}))\subset \mathrm{A}\)，从而 \(\widehat{u} (\overline{j(\mathrm{V})})\subset \mathrm{A}\)。而 \(\overline{j(\mathrm{V})}\) 是 \(\widehat{\mathrm{E}}\) 中 0 的一个邻域（TG, III, p. 24, prop. 7）。因此，\(\widehat{u}\) 的像包含于 F，且连续线性映射 \(v\colon \widehat{\mathrm{E}}\rightarrow \mathrm{F}\)（由 \(\widehat{u}\) 过渡到子空间得到的）是紧的。由于 \(j(\mathrm{E})\) 在 \(\widehat{\mathrm{E}}\) 中稠密，映射 \(v\) 是从 \(\widehat{\mathrm{E}}\) 到 F 的满足 \(u = v\circ j\) 的唯一连续线性映射。

命题 5. --- 设 E 是局部凸空间，F 是分离的拓扑向量空间，u 是从 E 到 F 的紧线性映射。则存在一个巴拿赫空间 G、一个从 E 到 G 的连续线性映射 v 和一个从 G 到 F 的紧线性映射 w，使得 \(u = w \circ v\)。

取 E 中 0 的一个邻域 V，使得 \(u(\mathrm{V})\) 在 F 中相对紧。设 \(p\) 是 E 上的一个连续半范数，使得 V 包含 \(p\) 的单位球（EVT, II, p. 26, cor.）。用 \(\mathrm{E}_p\) 表示给 E 赋以半范数 \(p\) 所得到的半赋范空间。映射 \(u\) 是从 \(\mathrm{E}_p\) 到 F 的紧线性映射。\(\mathrm{E}_p\) 的分离完备化 G 是一个巴拿赫空间。用 \(v\) 表示从 \(\mathrm{E}_p\) 到 G 的典范线性映射。由命题 4，存在一个紧线性映射 \(w: \mathrm{G} \to \mathrm{F}\)，使得 \(u = w \circ v\)。另一方面，v 是从 E 到 G 的连续映射，因为 E 的拓扑比 \(\mathrm{E}_p\) 的拓扑更细。

